%% ---------------------------------------------------------------------------
%% Frozen foundation encoders for re-identification — IEEE TCSVT submission
%%
%% Build:  latexmk -pdf main.tex        (after: python tools/gen_tables.py)
%%
%% Every number in this document comes from generated/, which is written by
%% tools/gen_tables.py out of results/runs. Do not type a metric into the prose.
%% ---------------------------------------------------------------------------
\documentclass[journal]{IEEEtran}

\usepackage[T1]{fontenc}
\usepackage[utf8]{inputenc}
\usepackage{amsmath,amssymb}
\usepackage{booktabs}
\usepackage{graphicx}
\usepackage{multirow}
\usepackage{xcolor}
\usepackage{url}
\usepackage[hidelinks]{hyperref}
\usepackage[capitalise,nameinlink]{cleveref}

% A number that went the wrong way. Named rather than coloured inline, so that the one
% decision about how a regression looks lives in one place.
\newcommand{\loss}[1]{\textcolor{black!65}{#1}}

% Section marker for work that is scaffolded but not yet measured. Delete the macro's body
% before submission and every such marker disappears with it.
\newcommand{\scaffold}[1]{\par\noindent\textit{[#1]}\par}

% generated by tools/gen_tables.py from results/runs — do not edit
\newcommand{\nRuns}{672}
\newcommand{\nEncoderConfigs}{12}
\newcommand{\nEncoders}{7}
\newcommand{\nDatasets}{7}
\newcommand{\nProtocols}{8}
\newcommand{\nHeads}{7}
\newcommand{\nHeadRecipes}{4}
\newcommand{\nFitSplits}{2}
\newcommand{\nHeadIdentities}{751}
\newcommand{\nHeadImages}{12,936}
\newcommand{\bestMarketMAP}{0.7181}
\newcommand{\bestMarketRone}{0.8872}
\newcommand{\leadMarketFrozen}{0.0610}
\newcommand{\leadMarketPCA}{0.0690}
\newcommand{\leadMarketLinear}{0.3998}
\newcommand{\leadMarketArc}{0.7087}
\newcommand{\leadMarketMINP}{0.3202}
\newcommand{\leadVraiFrozen}{0.1739}
\newcommand{\leadVraiPCA}{0.1625}
\newcommand{\leadVraiArc}{0.2276}
\newcommand{\leadOccludedFrozen}{0.4560}
\newcommand{\leadOccludedPCA}{0.4563}
\newcommand{\leadOccludedArc}{0.6459}
\newcommand{\teacherMarketMINP}{0.2038}
\newcommand{\leadMarketGain}{11.6}
\newcommand{\minpGap}{1.57}
\newcommand{\rankSiglipgFrozen}{1}
\newcommand{\rankSiglipgProbed}{6}
\newcommand{\mapSiglipgFrozen}{0.1051}
\newcommand{\mapSiglipgProbed}{0.6077}
\newcommand{\gainSiglipg}{5.8}
\newcommand{\rankSiglipsoFrozen}{2}
\newcommand{\rankSiglipsoProbed}{10}
\newcommand{\mapSiglipsoFrozen}{0.0772}
\newcommand{\mapSiglipsoProbed}{0.5137}
\newcommand{\gainSiglipso}{6.7}
\newcommand{\rankDinohpFrozen}{8}
\newcommand{\rankDinohpProbed}{5}
\newcommand{\mapDinohpFrozen}{0.0487}
\newcommand{\mapDinohpProbed}{0.6548}
\newcommand{\gainDinohp}{13.5}
\newcommand{\rankRadiohFrozen}{5}
\newcommand{\rankRadiohProbed}{3}
\newcommand{\mapRadiohFrozen}{0.0610}
\newcommand{\mapRadiohProbed}{0.7087}
\newcommand{\gainRadioh}{11.6}
\newcommand{\studentMarket}{0.7087}
\newcommand{\teacherMarket}{0.6077}
\newcommand{\dinoMarket}{0.6548}
\newcommand{\siglipsmMarket}{0.5137}
\newcommand{\clipMarket}{0.2887}
\newcommand{\clipcropMarket}{0.1744}
\newcommand{\studentCuhk}{0.1528}
\newcommand{\teacherCuhk}{0.2826}
\newcommand{\dinoCuhk}{0.0971}
\newcommand{\siglipsmCuhk}{0.1743}
\newcommand{\clipCuhk}{0.0165}
\newcommand{\clipcropCuhk}{0.0183}
\newcommand{\studentOccluded}{0.6459}
\newcommand{\teacherOccluded}{0.6494}
\newcommand{\dinoOccluded}{0.3560}
\newcommand{\siglipsmOccluded}{0.5576}
\newcommand{\clipOccluded}{0.4137}
\newcommand{\clipcropOccluded}{0.4370}
\newcommand{\studentVrai}{0.2276}
\newcommand{\teacherVrai}{0.2290}
\newcommand{\dinoVrai}{0.2467}
\newcommand{\siglipsmVrai}{0.1772}
\newcommand{\clipVrai}{0.0394}
\newcommand{\clipcropVrai}{0.0346}
\newcommand{\studentVric}{0.1909}
\newcommand{\teacherVric}{0.1148}
\newcommand{\dinoVric}{0.3588}
\newcommand{\siglipsmVric}{0.0746}
\newcommand{\clipVric}{0.0349}
\newcommand{\clipcropVric}{0.0489}
\newcommand{\studentCcvid}{0.6093}
\newcommand{\teacherCcvid}{0.6257}
\newcommand{\dinoCcvid}{0.4615}
\newcommand{\siglipsmCcvid}{0.5888}
\newcommand{\clipCcvid}{0.7947}
\newcommand{\clipcropCcvid}{0.6301}
\newcommand{\studentMsmt}{0.1645}
\newcommand{\teacherMsmt}{0.2307}
\newcommand{\dinoMsmt}{0.0669}
\newcommand{\siglipsmMsmt}{0.1999}
\newcommand{\clipMsmt}{0.0269}
\newcommand{\clipcropMsmt}{0.0160}
\newcommand{\spreadMarket}{1.08}
\newcommand{\deficitMarket}{+8}
\newcommand{\familyStudentMarket}{0.7087}
\newcommand{\familySiglipMarket}{0.6077}
\newcommand{\familyDinoMarket}{0.6548}
\newcommand{\spreadCuhk}{2.91}
\newcommand{\deficitCuhk}{-46}
\newcommand{\familyStudentCuhk}{0.1528}
\newcommand{\familySiglipCuhk}{0.2826}
\newcommand{\familyDinoCuhk}{0.0971}
\newcommand{\spreadOccluded}{1.82}
\newcommand{\deficitOccluded}{+0}
\newcommand{\familyStudentOccluded}{0.6510}
\newcommand{\familySiglipOccluded}{0.6494}
\newcommand{\familyDinoOccluded}{0.3560}
\newcommand{\spreadVrai}{1.08}
\newcommand{\deficitVrai}{-8}
\newcommand{\familyStudentVrai}{0.2276}
\newcommand{\familySiglipVrai}{0.2290}
\newcommand{\familyDinoVrai}{0.2467}
\newcommand{\spreadVric}{3.12}
\newcommand{\deficitVric}{-45}
\newcommand{\familyStudentVric}{0.1971}
\newcommand{\familySiglipVric}{0.1148}
\newcommand{\familyDinoVric}{0.3588}
\newcommand{\spreadCcvid}{1.36}
\newcommand{\deficitCcvid}{-0}
\newcommand{\familyStudentCcvid}{0.6254}
\newcommand{\familySiglipCcvid}{0.6257}
\newcommand{\familyDinoCcvid}{0.4615}
\newcommand{\spreadMsmt}{3.35}
\newcommand{\deficitMsmt}{-29}
\newcommand{\familyStudentMsmt}{0.1645}
\newcommand{\familySiglipMsmt}{0.2307}
\newcommand{\familyDinoMsmt}{0.0689}
\newcommand{\fitFrozenMarket}{0.056}
\newcommand{\fitPcaMarketHere}{0.064}
\newcommand{\fitPcaMarketElse}{0.058}
\newcommand{\fitSharePcaMarket}{26\%}
\newcommand{\fitSwingPcaMarket}{0.006}
\newcommand{\fitLinearMarketHere}{0.304}
\newcommand{\fitLinearMarketElse}{0.183}
\newcommand{\fitShareLinearMarket}{51\%}
\newcommand{\fitSwingLinearMarket}{0.121}
\newcommand{\fitArcfaceMarketHere}{0.568}
\newcommand{\fitArcfaceMarketElse}{0.292}
\newcommand{\fitShareArcfaceMarket}{46\%}
\newcommand{\fitSwingArcfaceMarket}{0.276}
\newcommand{\fitFrozenMsmt}{0.028}
\newcommand{\fitPcaMsmtHere}{0.033}
\newcommand{\fitPcaMsmtElse}{0.033}
\newcommand{\fitSharePcaMsmt}{92\%}
\newcommand{\fitSwingPcaMsmt}{0.000}
\newcommand{\fitLinearMsmtHere}{0.206}
\newcommand{\fitLinearMsmtElse}{0.082}
\newcommand{\fitShareLinearMsmt}{31\%}
\newcommand{\fitSwingLinearMsmt}{0.124}
\newcommand{\fitArcfaceMsmtHere}{0.399}
\newcommand{\fitArcfaceMsmtElse}{0.129}
\newcommand{\fitShareArcfaceMsmt}{27\%}
\newcommand{\fitSwingArcfaceMsmt}{0.269}
\newcommand{\nHeadImagesMsmt}{30,248}
\newcommand{\nHeadIdentitiesMsmt}{1,041}


\begin{document}

\title{The Readout Decides the Winner: Probing Frozen\\Foundation Encoders for Person and
Vehicle Re-Identification}

\author{Karol~Majek% <-- affiliation, ORCID, funding and co-authors before submission
\thanks{Manuscript prepared \today. Code, run records and the exact protocol digests behind
every number in this paper are released with the \texttt{reidbench} package (MIT).}}

\markboth{IEEE Transactions on Circuits and Systems for Video Technology}%
{Majek: The Readout Decides the Winner}

\maketitle

\begin{abstract}
Frozen vision foundation models are increasingly proposed as general-purpose encoders for
object re-identification (ReID), and are compared to one another by cosine similarity over
their raw features. We show that this comparison does not measure what it is taken to
measure. Across \nEncoderConfigs{} encoder-resolution configurations of \nEncoders{} modern
backbones --- agglomerative (C-RADIOv4), language-aligned (SigLIP2, CLIP) and self-supervised
(DINOv3) --- evaluated under \nProtocols{} exact retrieval protocols over \nDatasets{} person
and vehicle datasets with \nHeadRecipes{} readout heads fitted on \nFitSplits{} training
splits, \nRuns{} run records in total, we find that
\emph{the choice of readout, not the choice of backbone, determines the ranking}. Fitting a
single 512-dimensional affine map on \nHeadImages{} images of \nHeadIdentities{}
identities moves SigLIP2-g from rank \rankSiglipgFrozen{} to rank \rankSiglipgProbed{} on
Market-1501 and DINOv3-H+ from rank \rankDinohpFrozen{} to rank \rankDinohpProbed{}, while
multiplying frozen mAP by between \gainSiglipg{}$\times$ and
\gainDinohp{}$\times$ depending on the backbone. Contrastively pre-trained encoders gain least
because their objective already optimised the cosine geometry that zero-shot evaluation reads;
frozen-cosine benchmarking therefore flatters them systematically. A matched PCA control that
sees the same images and none of the labels isolates the effect: the gain is the supervision,
not the bottleneck. Refitting every head on a second training split then splits that
supervision in two. A head that has never seen a target label still carries
\fitShareArcfaceMarket{} of the gain, so most of what a readout supplies is a transferable
metric geometry; the remainder is domain fitting, and it is worth \fitSwingArcfaceMarket{} mAP
on Market-1501 --- more than the gap between several of the backbones compared here. We then use the same substrate to ask whether agglomerative distillation
costs the instance-level margin ReID depends on. It does not --- in-domain the student beats
its own teacher by a wider factor on mINP (\minpGap{}$\times$) than on mAP --- but its
advantage does not survive the move off the head's training domain, where it regresses toward
the mean of teachers that disagree sharply by domain. Two nuisance variables, the resize
geometry used to fit a crop to the encoder's input and the input resolution itself, are each
worth more than the gap between several of the backbones compared here. All protocols, run
records, provenance digests and licence records are released as an MIT-licensed package, and
the tables in this paper are generated from those records rather than transcribed.
\end{abstract}


\begin{IEEEkeywords}
Re-identification, foundation models, frozen features, linear probing, knowledge distillation,
evaluation protocols, vehicle re-identification, reproducibility.
\end{IEEEkeywords}

\section{Introduction}
\IEEEPARstart{O}{bject} re-identification asks whether two images, taken by different cameras
at different times, show the same individual. A decade of supervised work on person
benchmarks has driven in-domain accuracy close to saturation \cite{ye2022deep,zhang2024review},
and the field's attention has moved to two questions that in-domain accuracy does not answer:
how a representation behaves off the domain it was trained on \cite{kim2025dgsurvey}, and
whether a general-purpose visual encoder can replace the specialist entirely
\cite{li2023clipreid,paradigm2026}.

The second question is now asked routinely, and it is usually asked in the same way. A
foundation backbone is frozen, its features are extracted, queries are ranked against a
gallery by cosine similarity, and the resulting mAP is reported beside another backbone's.
The comparison is cheap, it needs no training, and it appears to isolate the representation
from everything else. \textbf{This paper's central claim is that it does not, and that the
ranking it produces is not stable under a change that a practitioner would regard as an
implementation detail.}

Concretely, we take \nEncoders{} modern backbones at \nEncoderConfigs{}
encoder-resolution configurations, spanning three pre-training paradigms --- agglomerative
distillation (C-RADIOv4 \cite{ranzinger2024amradio,heinrich2025radiov25}), language-image
contrastive learning (SigLIP2 \cite{tschannen2025siglip2}, CLIP \cite{radford2021clip}) and
self-supervised distillation (DINOv3 \cite{simeoni2025dinov3}) --- and evaluate each one under
four readouts over the identical cached features: the frozen encoder itself, a label-free PCA
projection, a linear head, and an ArcFace \cite{deng2019arcface} head. All three fitted heads
see one training split and one only, \texttt{market1501/train}. We evaluate every combination
on \nDatasets{} datasets under \nProtocols{} exact protocols, covering person and vehicle
re-identification, occlusion, cloth change, degraded detector boxes and aerial viewpoints.

Three findings follow, and the first is about methodology rather than about any backbone.

\subsection*{The readout decides the ranking}
On Market-1501, SigLIP2-g is the best encoder in this study when read frozen
(\mapSiglipgFrozen{} mAP, rank \rankSiglipgFrozen{}) and the \rankSiglipgProbed{}th when the
same weights are read through one affine map fitted on \nHeadImages{} images
(\mapSiglipgProbed{}). DINOv3-H+ moves the other way, from rank \rankDinohpFrozen{} to rank
\rankDinohpProbed{}. Nothing about the weights, the images, the protocol or the metric changed
between those two rankings --- only how the features were read. The mechanism is not
mysterious: a contrastive objective optimises cosine geometry directly, so a
contrastively pre-trained encoder is already read out near-optimally and an affine map has
little headroom left (\gainSiglipg{}$\times$), whereas encoders trained on other objectives
carry the information in a geometry cosine similarity cannot see, and a head recovers it
(\gainDinohp{}$\times$). Zero-shot frozen-cosine benchmarking is therefore not a neutral
instrument; it measures readout alignment as much as representation quality, and it flatters
one pre-training family in particular.

\subsection*{Agglomeration keeps the margin, in-domain}
Agglomerative backbones distil several teachers into one student
\cite{ranzinger2024amradio,heinrich2025radiov25}, and a standing concern is that averaging
teacher objectives erodes the fine-grained instance separability that re-identification, unlike
classification, depends on. Because C-RADIOv4's actual teachers are in our encoder set, we can
test this directly rather than by analogy. In-domain the concern is not borne out and the
opposite holds on exactly the metric it is about: against its own teacher SigLIP2-g, the
student's advantage is \minpGap{}$\times$ on mINP --- the rank of the \emph{hardest} true match
--- against a smaller factor on mAP. Off the head's training domain the advantage disappears,
and the shape of its disappearance is regular: the student lands between its two teacher
families wherever those families disagree, which they do sharply and by domain.

\subsection*{Two nuisance variables outweigh the model choice}
The transform that fits a $64\times128$ pedestrian crop to a $224\times224$ input is not a
detail. \texttt{timm}'s default evaluation transform short-side-resizes and centre-crops,
which on such a crop discards the head and the feet; the alternative is to resize to the target
directly. Measured on identical CLIP weights with geometry as the only variable, the choice is
worth a large fraction of the probed Market-1501 score. Input resolution is worth a comparable
amount in the other direction. Either variable, left unpinned, is large enough to reverse the
ordering of several backbone pairs in this study --- which is a reproducibility problem for
frozen-encoder comparisons generally, not a property of these particular models.

\subsection*{Contributions}
\begin{enumerate}
  \item A controlled demonstration that readout choice reorders frozen-encoder rankings on a
        standard ReID benchmark, with a matched label-free control that attributes the effect
        to supervision rather than to dimensionality (\cref{sec:readout,sec:heads}).
  \item The first evaluation of agglomerative vision foundation models as ReID encoders,
        together with an ablation against their own teachers at matched capacity, separating
        an in-domain gain from a transfer cost with a predictable shape
        (\cref{sec:ablation}).
  \item Quantification of two preprocessing variables --- resize geometry and input resolution
        --- that are routinely left unreported and that are individually larger than the
        differences such comparisons are used to establish (\cref{sec:geometry}).
  \item An MIT-licensed evaluation substrate in which a protocol is a value with a digest, a
        run record carries its own provenance, and every table in this paper is generated from
        those records rather than transcribed (\cref{sec:substrate}).
\end{enumerate}

We deliberately train no backbone. Every number here is one forward pass over cached crops
plus, at most, under a minute of head fitting on a single consumer GPU, which is what makes a
grid of this size affordable and what makes each cell exactly reproducible from its digest.

\section{Related Work}\label{sec:related}

\subsection{Re-identification with frozen and foundation encoders}
CLIP-ReID \cite{li2023clipreid} established that a language-image backbone is a strong
initialisation for person ReID, and a family of adaptation methods followed --- prompt
learning, parameter-efficient tuning and instruction conditioning
\cite{he2024mambapro,chen2024instructreid,zhu2022pass}. These works fine-tune. The
complementary question, what a foundation encoder is worth \emph{without} adaptation, is
answered less often and less consistently. A 2026 cross-paradigm study
\cite{paradigm2026} compares supervised, self-supervised and language-aligned models over 11
backbones and 9 datasets and is the closest prior work to ours in scope; a parallel line
evaluates multimodal LLMs as the matcher itself \cite{mmreid2025}. Neither varies the readout
as an experimental axis: each fixes a single evaluation convention and reports backbone
differences under it. Our contribution is orthogonal and, we argue, prior to theirs --- if the
ranking depends on the convention, then a table that fixes one convention is reporting a
property of the pair, not of the backbones.

\subsection{Agglomerative vision foundation models}
AM-RADIO \cite{ranzinger2024amradio} showed that several heterogeneous teachers --- a
language-aligned model, a self-supervised model and a segmentation model --- can be distilled
into a single student that matches or exceeds each teacher on that teacher's own benchmarks,
and RADIOv2.5 \cite{heinrich2025radiov25} addressed the mode-switching and resolution
pathologies that the first version exhibited. The line has since been extended by other
agglomerative students \cite{eupe2026} and by mixture-of-experts variants
\cite{amoe2025}. What this family has not been evaluated on is instance-level retrieval, where
the relevant quantity is not whether a class is separable but whether one individual is
separable from a visually near-identical other. Distillation from an averaged teacher signal
is a plausible threat to exactly that margin, and it is the hypothesis \cref{sec:ablation}
tests against C-RADIOv4's own teachers.

\subsection{Probing a frozen representation}
Linear probing is the standard instrument for comparing frozen representations, inherited from
the self-supervised classification literature \cite{chen2020simclr,caron2021dino,oquab2024dinov2}.
Two conventions travel under the name and are not interchangeable. The classification
convention trains a linear classifier on frozen features and reports \emph{the classifier's}
accuracy; the retrieval convention discards the classifier and reports mAP over the features
it reshaped. Foundation-model reports overwhelmingly use the first, ReID uses the second, and
numbers from one are quoted beside numbers from the other. Within ReID, the readout that
matters in practice is a margin-based one --- ArcFace \cite{deng2019arcface},
CosFace \cite{wang2018cosface} or Circle loss \cite{sun2020circle} --- because it optimises the
angular geometry the retrieval metric reads, which is precisely why substituting it for a plain
linear head is not a neutral implementation choice. We report the retrieval convention
throughout and treat the head as an axis of the experiment rather than as part of its setup.

\subsection{Evaluation protocol and its failure modes}
Market-1501 \cite{zheng2015scalable} fixed the query--gallery--junk convention that most person
benchmarks still follow; MSMT17 \cite{wei2018msmt} scaled it; CUHK03-NP
\cite{zhong2017reranking} supplied the detector-box regime and the 767/700 split that made
CUHK03 comparable across papers. mINP \cite{ye2022deep} was introduced precisely because mAP is
dominated by easy true matches and says little about the hardest one, which is the quantity a
deployed system fails on. Beyond ReID, OpenOOD \cite{zhang2023openood} is the clearest recent
demonstration that protocol discipline --- fixed splits, tuning on validation data only,
stratified difficulty --- changes published conclusions in a whole subfield, and the
open-set ReID literature \cite{liao2014operid,li2018adversarial,wang2020gom} has repeatedly
observed that closed-set rank metrics assume the answer is present in the gallery. Our
substrate (\cref{sec:substrate}) is built on the same premise: an evaluation protocol should
be a value with an identity, so that two numbers can be checked for comparability rather than
assumed to be comparable.

\subsection{Benchmark toolboxes}
Torchreid \cite{zhou2019torchreid} and FastReID \cite{he2023fastreid} are the field's reference
implementations and are training frameworks first; their evaluation code is a component of a
training loop rather than an independently usable artifact, and neither exposes an
agglomerative backbone, a frozen-probe harness, or an open-set protocol. We release the
evaluation half as a standalone package with no training loop, no dataset mirror and no torch
dependency in its core, which is what allows a run record to be re-scored on a machine that
cannot run the encoder that produced it.

\section{The Evaluation Substrate}\label{sec:substrate}
A study of this shape --- \nRuns{} run records over four axes --- is only as trustworthy as its
bookkeeping. Two rows in a results table are a comparison if and only if they were produced
over the same images, under the same rules, and differ in the one variable being claimed. That
condition is easy to state and easy to violate silently. This section describes the substrate
we built to make the condition checkable, because several of the results in this paper were
originally a different number, and were corrected by a check the substrate made possible.

\subsection{Two boolean matrices, and nothing downstream of them}
The substrate reduces every re-identification protocol to a pair of boolean matrices over a
selected query and gallery set:
\begin{align}
  \mathrm{rel}[q,g]   &= \mathbb{1}[\text{$q$ and $g$ are the same identity}], \\
  \mathrm{valid}[q,g] &= \mathbb{1}[\text{this pair is scored}].
\end{align}
Four functions close over them:
\begin{align}
  \textsf{select}&: (\text{manifest}, \text{protocol}) \rightarrow (q, g, \mathrm{rel},
     \mathrm{valid}) \nonumber\\
  \textsf{score}&: (X_q, X_g) \rightarrow S \nonumber\\
  \textsf{measure}&: (S, \mathrm{rel}, \mathrm{valid}) \rightarrow
     (\text{metrics}, \text{curves}) \nonumber\\
  \textsf{describe}&: \text{value} \rightarrow \text{canonical, hashable record.} \nonumber
\end{align}
Once $\mathrm{rel}$ and $\mathrm{valid}$ exist, cameras, junk identities, tracklets, clothing
labels, modality and open-set non-mates have all disappeared: mAP, CMC, mINP and the open-set
measures are functions of $(S, \mathrm{rel}, \mathrm{valid})$ and nothing below
\textsf{select} has heard of re-identification. The same-camera exclusion that every
Market-1501 number depends on is an entry in $\mathrm{valid}$, not a branch inside a metric,
which is why a dataset that labels no cameras --- Occluded-REID does not --- is expressible
rather than a special case.

\subsection{A protocol is a value with a digest}
A protocol is a named, versioned value naming its query split, its gallery split and its
exclusion rules, and it carries a content digest. Every run record stores that digest beside
the digest of the manifest it selected over. Two rows are comparable when both digests match,
and the substrate reports it rather than leaving it to the reader. This is what allows
CCVID's two protocols --- with and without the cloth-changing restriction --- to be reported
side by side without one silently being read as the other's improvement: they are two
different questions with two different digests, and a table that pooled them would emphasise
the easier one and call it the better system.

\subsection{Provenance, and the licence question}
Every run record carries the encoder specification, the content-addressed feature-cache key,
library versions, GPU and driver, and the source revision with a dirty-tree flag. Model and
dataset licences are indexed in the same package and reported with the metrics rather than
maintained beside them: of the checkpoints evaluated here, only the two C-RADIOv4 weights
carry a verified licence permitting commercial use, and the substrate says so on every run.
For a paper this is a footnote; for anyone acting on the paper's recommendation it is the
first question.

\subsection{Guards, and what they caught}
Before the first number from a new dataset adapter is believed, we score two synthetic feature
sets under the adapter's protocol: \emph{perfect} features, in which the relevance matrix is
used as its own similarity, and \emph{random} ones. A correct adapter returns mAP $1.0$ and a
value near the protocol's chance floor respectively. On CUHK03-NP this gave $1.0000$ and
$0.0022$, against Market-1501's $1.0000$ and $0.0015$, which is what distinguishes a hard
dataset from a broken adapter --- a distinction no amount of staring at a low score can make.

The substrate also caught the largest error in this study. The teacher backbones are served
through \texttt{timm}, whose default evaluation transform short-side-resizes and centre-crops;
the \texttt{torchhub} path used for C-RADIOv4 resizes to the target directly. On a
$64\times128$ pedestrian crop the former retains roughly the middle 45\% of the
body. Run as first configured, the teacher ablation would have measured a preprocessing
difference and reported it as a distillation result, with the same confidence and the opposite
conclusion. The fix was to make the resize geometry a hashed field of the encoder
specification, so that the two geometries occupy different cache keys and appear as different
rows (\cref{sec:geometry}) rather than as one contradictory one.

\section{Experimental Design}\label{sec:setup}
The experiment is a cross product of four axes --- encoder, input geometry, readout head and
dataset --- evaluated under exact protocols. \Cref{tab:encoders} lists the encoders,
\cref{tab:datasets} the datasets and protocols, and \cref{tab:coverage} which cells of the
product are measured, so that an absence is visible as an absence rather than as a blank.

\begin{table*}[t]
\centering
\caption{The encoders. Every checkpoint is frozen: no gradient reaches a backbone anywhere in
this study. \emph{fit} is how an image is made to match the input size --- \texttt{squash}
resizes to the target, \texttt{crop} short-side-resizes and centre-crops. C-RADIOv4 distils
SigLIP2-g, DINOv3 and a segmentation teacher, so rows 1--6 are students of rows 7--10's
families.}
\label{tab:encoders}
\small
\setlength{\tabcolsep}{4pt}
% generated by tools/gen_tables.py from results/runs — edit the runs, not this
\begin{tabular}{lllllp{0.26\linewidth}}
\toprule
encoder & pre-training & params (M) & input & fit & checkpoint \\
\midrule
C-RADIOv4-H & agglomerative & 653 & 224$\times$224 & squash & \texttt{\small torchhub:\allowbreak NVlabs/\allowbreak RADIO/\allowbreak c-radio\_\allowbreak v4-h} \\
C-RADIOv4-H & agglomerative & 653 & 256$\times$128 & squash & \texttt{\small torchhub:\allowbreak NVlabs/\allowbreak RADIO/\allowbreak c-radio\_\allowbreak v4-h} \\
C-RADIOv4-H & agglomerative & 653 & \textit{native} & squash & \texttt{\small torchhub:\allowbreak NVlabs/\allowbreak RADIO/\allowbreak c-radio\_\allowbreak v4-h} \\
C-RADIOv4-SO400M & agglomerative & 431 & 224$\times$224 & squash & \texttt{\small torchhub:\allowbreak NVlabs/\allowbreak RADIO/\allowbreak c-radio\_\allowbreak v4-so400m} \\
C-RADIOv4-SO400M & agglomerative & 431 & 256$\times$128 & squash & \texttt{\small torchhub:\allowbreak NVlabs/\allowbreak RADIO/\allowbreak c-radio\_\allowbreak v4-so400m} \\
C-RADIOv4-SO400M & agglomerative & 431 & \textit{native} & squash & \texttt{\small torchhub:\allowbreak NVlabs/\allowbreak RADIO/\allowbreak c-radio\_\allowbreak v4-so400m} \\
SigLIP2-g & language-aligned & 1163 & 256$\times$256 & squash & \texttt{\small timm:\allowbreak vit\_\allowbreak giantopt\_\allowbreak patch16\_\allowbreak siglip\_\allowbreak 256.v2\_\allowbreak webli} \\
SigLIP2-SO400M & language-aligned & 428 & 224$\times$224 & squash & \texttt{\small timm:\allowbreak vit\_\allowbreak so400m\_\allowbreak patch14\_\allowbreak siglip\_\allowbreak 224.v2\_\allowbreak webli} \\
DINOv3-H+ & self-supervised & 840 & 224$\times$224 & squash & \texttt{\small timm:\allowbreak vit\_\allowbreak huge\_\allowbreak plus\_\allowbreak patch16\_\allowbreak dinov3.lvd1689m} \\
DINOv3-L & self-supervised & 303 & 224$\times$224 & squash & \texttt{\small timm:\allowbreak vit\_\allowbreak large\_\allowbreak patch16\_\allowbreak dinov3.lvd1689m} \\
CLIP-B/16 & language-aligned & 87 & 224$\times$224 & squash & \texttt{\small timm:\allowbreak vit\_\allowbreak base\_\allowbreak patch16\_\allowbreak clip\_\allowbreak 224.openai} \\
CLIP-B/16 & language-aligned & 87 & 224$\times$224 & crop & \texttt{\small timm:\allowbreak vit\_\allowbreak base\_\allowbreak patch16\_\allowbreak clip\_\allowbreak 224.openai} \\
\bottomrule
\end{tabular}

\end{table*}

\subsection{Encoders}
\nEncoders{} backbones at \nEncoderConfigs{} encoder-resolution configurations span three
pre-training paradigms. \textbf{C-RADIOv4} (H, 653M; SO400M, 431M) is
agglomerative: one student distilled from SigLIP2-g, DINOv3-7B and a segmentation teacher
\cite{ranzinger2024amradio,heinrich2025radiov25}. \textbf{SigLIP2}
\cite{tschannen2025siglip2} appears at 1163M --- the literal teacher --- and at
428M, capacity-matched to C-RADIOv4-SO400M. \textbf{DINOv3}
\cite{simeoni2025dinov3} appears at 840M and 303M, bracketing C-RADIOv4-H's
653M from above and below; note that the true teacher is DINOv3-7B, which does not fit
on the card this study ran on, so the DINOv3 rows stand in for a teacher family rather than
being the teacher. \textbf{CLIP ViT-B/16} \cite{radford2021clip} is the small, well-understood
control, and is the one encoder present under both resize geometries.

Features are the encoder's pooled summary representation at \texttt{fp32}, cached to a
content-addressed store keyed on the encoder description and the manifest digest. C-RADIOv4 is
additionally run at $256\times128$ and at each image's native size, snapped to the model's
patch grid and capped at 512 pixels, which is the resolution axis of \cref{sec:geometry}.

\subsection{Readout heads}
A \emph{head} is a small map fitted on the frozen encoder's cached features and nothing else.
\nHeadRecipes{} recipes are reported, and \texttt{none} is one of them:
\begin{itemize}
  \item \textbf{none} --- the frozen encoder, ranked by cosine similarity. The zero-shot
        convention, and the row every other row is read against.
  \item \textbf{pca} --- a 512-dimensional PCA projection fitted on
        \texttt{market1501/train}. It sees the same \nHeadImages{} images as the supervised
        heads and none of the \nHeadIdentities{} identity labels. This is the control that
        separates \emph{supervision} from \emph{dimensionality}.
  \item \textbf{linear} --- a 512-dimensional affine map trained with softmax
        cross-entropy over the \nHeadIdentities{} training identities; the classifier is
        discarded and retrieval is scored on the projected features.
  \item \textbf{arcface} --- the same map trained with an additive angular margin
        \cite{deng2019arcface} ($m=0.3$, $s=30$).
\end{itemize}
All three fitted heads share one budget --- 100 epochs, Adam, lr $10^{-3}$, weight decay
$5\times10^{-4}$, batch 256, $\ell_2$-normalised inputs, seed 0 --- fixed on training accuracy
before any retrieval number was read. Fitting takes well under a minute per encoder on the
study's single GPU. Except where \cref{tab:fitsplit} says otherwise, every head is fitted on
\texttt{market1501/train} and applied unchanged everywhere, which makes exactly one column of
every results table in-domain supervised and every other column cross-domain transfer with no
target label ever seen. That asymmetry is the most important single fact about these tables, is
why the regime is named in \cref{tab:datasets}, and is measured rather than assumed in
\cref{sec:heads}, where the same recipes are refitted on a second training split.

\begin{table*}[t]
\centering
\caption{Datasets and protocols. \emph{regime} is relative to the heads, all of which are
fitted on \texttt{market1501/train}; only Market-1501 is in-domain. Query and gallery sizes are
read from the run records, so a cross-dataset comparison of absolute scores can be checked
against the search size that produced it.}
\label{tab:datasets}
\scriptsize
\setlength{\tabcolsep}{3pt}
% generated by tools/gen_tables.py from results/runs — edit the runs, not this
\begin{tabular}{llllrr}
\toprule
dataset & object & regime & protocol & queries & gallery \\
\midrule
Market-1501 & person & in-domain & \texttt{market1501/official@1} & 3,368 & 19,732 \\
MSMT17 & person & transfer & \texttt{msmt17/official@1} & 11,659 & 82,161 \\
CUHK03-NP & person & transfer & \texttt{cuhk03/detected-767@1} & 1,400 & 5,332 \\
Occluded-REID & person & transfer & \texttt{occluded-reid/occluded-vs-whole@1} & 1,000 & 1,000 \\
CCVID & person & transfer & \texttt{ccvid/tracklet-cloth-changing@1} & 834 & 1,074 \\
CCVID & person & transfer & \texttt{ccvid/tracklet@1} & 834 & 1,074 \\
VRIC & vehicle & transfer & \texttt{vric/official@1} & 2,811 & 2,811 \\
VRAI & vehicle & transfer & \texttt{vrai/train-cross-camera@1} & 6,302 & 32,338 \\
\bottomrule
\end{tabular}

\end{table*}

\subsection{Datasets and protocols}
\nDatasets{} datasets under \nProtocols{} protocols cover five regimes of difficulty.
\textbf{Market-1501} \cite{zheng2015scalable} is the source domain and the only in-domain
column. \textbf{MSMT17} \cite{wei2018msmt} is the large-scale person transfer target.
\textbf{CUHK03-NP} \cite{zhong2017reranking}, detected boxes under the 767/700 split, is the
degraded-detector regime and the one closest to what a real detector hands a ReID model.
\textbf{Occluded-REID} \cite{zhuo2018occluded} ships no standard split and labels no cameras,
so the whole set is used, occluded probes against whole-body gallery, and the same-camera
exclusion every Market number depends on does not exist there. \textbf{CCVID}
\cite{gu2022ccvid} is video and cloth-changing; frames are thinned to eight per tracklet before
encoding and mean-pooled into one vector per tracklet afterwards, and both protocols --- with
and without the cloth-changing restriction --- are reported. \textbf{VRIC} \cite{kanaci2018vric}
and \textbf{VRAI} \cite{wang2019vrai} are vehicles, the latter aerial. VRAI's test identities
are withheld by its release, so the only protocol that scores offline runs over its training
split, querying the first frame of each camera-1 trajectory against every camera-2 training
image; that is a legitimate zero-shot number for an encoder that never saw VRAI, and it is
sound for the head rows here \emph{only} because no head in this study has ever seen a VRAI
image.

We do not use DukeMTMC or any dataset derived from it. The exclusion is enforced in code with
no override.

\begin{table*}[t]
\centering
\caption{Coverage: run records per (encoder, dataset) pair, four per cell being the four heads
and eight where the dataset ships two protocols. The MSMT17 column is the study's one
systematic gap and is discussed in \cref{sec:limitations}.}
\label{tab:coverage}
\scriptsize
\setlength{\tabcolsep}{3pt}
% generated by tools/gen_tables.py from results/runs — edit the runs, not this
\begin{tabular}{llrrrrrrr}
\toprule
encoder & input & Market & MSMT17 & CUHK03 & Occ-REID & CCVID & VRIC & VRAI \\
\midrule
C-RADIOv4-H & 224$\times$224 & 7 & 7 & 7 & 7 & 7 & 7 & 7 \\
C-RADIOv4-H & 256$\times$128 & 7 & 7 & 7 & 7 & 7 & 7 & 7 \\
C-RADIOv4-H & \textit{native} & 7 & 7 & 7 & 7 & 7 & 7 & 7 \\
C-RADIOv4-SO400M & 224$\times$224 & 7 & 7 & 7 & 7 & 7 & 7 & 7 \\
C-RADIOv4-SO400M & 256$\times$128 & 7 & 7 & 7 & 7 & 7 & 7 & 7 \\
C-RADIOv4-SO400M & \textit{native} & 7 & 7 & 7 & 7 & 7 & 7 & 7 \\
SigLIP2-g & 256$\times$256 & 7 & 7 & 7 & 7 & 7 & 7 & 7 \\
SigLIP2-SO400M & 224$\times$224 & 7 & 7 & 7 & 7 & 7 & 7 & 7 \\
DINOv3-H+ & 224$\times$224 & 7 & 7 & 7 & 7 & 7 & 7 & 7 \\
DINOv3-L & 224$\times$224 & 7 & 7 & 7 & 7 & 7 & 7 & 7 \\
CLIP-B/16 (squash) & 224$\times$224 & 7 & 7 & 7 & 7 & 7 & 7 & 7 \\
CLIP-B/16 (crop) & 224$\times$224 & 7 & 7 & 7 & 7 & 7 & 7 & 7 \\
\bottomrule
\end{tabular}

\end{table*}

\subsection{Metrics}
We report mAP and Rank-1 throughout, and mINP \cite{ye2022deep} where the claim is about the
hardest true match rather than the average one. Cross-dataset comparison of absolute values is
not meaningful here and we do not make it: Occluded-REID searches roughly a thousand images,
Market-1501 searches nearly twenty thousand, and VRAI over thirty thousand, so a gallery an
order of magnitude larger accounts for much of the difference between those columns before any
question of difficulty arises. Every claim in this paper is a within-column comparison, where
all rows share a protocol digest and a manifest digest.

\subsection{Noise}
Refitting the lead encoder's ArcFace head at three seeds and rescoring gives a standard
deviation of $0.0002$ mAP on Market-1501, $0.0050$ on Occluded-REID --- the noisiest set here,
with its thousand queries --- and $0.0009$ on VRAI. We treat differences below roughly twice
those values as ties and say so where it matters; several comparisons in \cref{sec:ablation}
turn on this.

\section{The Readout Decides the Ranking}\label{sec:readout}
\Cref{tab:readout} is the experiment this paper is built on. It holds the weights, the images,
the protocol and the metric fixed, and varies only how the cached features are read.

\begin{table*}[t]
\centering
\caption{Market-1501 mAP under each readout, over identical cached features. \emph{rank} is the
encoder's position among the \nEncoderConfigs{} configurations when read frozen and when read
through the ArcFace head; \emph{gain} is the ratio of the two. The two CLIP rows are the same
checkpoint at the same input size under the two resize geometries. Values at three decimals;
the text quotes four.}
\label{tab:readout}
% generated by tools/gen_tables.py from results/runs — edit the runs, not this
\begin{tabular}{llrrrrcr}
\toprule
encoder & input & frozen & PCA & linear & ArcFace & rank & gain \\
\midrule
C-RADIOv4-H & 224$\times$224 & 0.061 & 0.069 & 0.400 & 0.709 & 5\,$\to$\,3 & 11.6$\times$ \\
C-RADIOv4-H & 256$\times$128 & 0.059 & 0.066 & \textbf{0.404} & \textbf{0.718} & 6\,$\to$\,1 & 12.1$\times$ \\
C-RADIOv4-H & \textit{native} & 0.044 & 0.049 & 0.268 & 0.559 & 10\,$\to$\,9 & 12.6$\times$ \\
C-RADIOv4-SO400M & 224$\times$224 & 0.063 & 0.070 & 0.368 & 0.696 & 4\,$=$\,4 & 11.1$\times$ \\
C-RADIOv4-SO400M & 256$\times$128 & 0.063 & 0.070 & 0.384 & 0.718 & 3\,$\to$\,2 & 11.4$\times$ \\
C-RADIOv4-SO400M & \textit{native} & 0.047 & 0.053 & 0.271 & 0.573 & 9\,$\to$\,8 & 12.1$\times$ \\
SigLIP2-g & 256$\times$256 & \textbf{0.105} & \textbf{0.126} & 0.402 & 0.608 & 1\,$\to$\,6 & 5.8$\times$ \\
SigLIP2-SO400M & 224$\times$224 & 0.077 & 0.092 & 0.289 & 0.514 & 2\,$\to$\,10 & 6.7$\times$ \\
DINOv3-H+ & 224$\times$224 & 0.049 & 0.056 & 0.319 & 0.655 & 8\,$\to$\,5 & 13.5$\times$ \\
DINOv3-L & 224$\times$224 & 0.051 & 0.055 & 0.302 & 0.607 & 7\,$=$\,7 & 11.9$\times$ \\
CLIP-B/16 (squash) & 224$\times$224 & 0.029 & 0.034 & 0.142 & 0.289 & 11\,$=$\,11 & 10.0$\times$ \\
CLIP-B/16 (crop) & 224$\times$224 & 0.023 & 0.027 & 0.095 & 0.174 & 12\,$=$\,12 & 7.7$\times$ \\
\bottomrule
\end{tabular}

\end{table*}

\subsection{The ranking is not stable}
Read frozen, the two SigLIP2 checkpoints are the best encoders in this study on Market-1501,
at ranks \rankSiglipgFrozen{} and \rankSiglipsoFrozen{}. Read through one affine map fitted on
\nHeadImages{} images, they are ranked \rankSiglipgProbed{} and \rankSiglipsoProbed{}. DINOv3-H+
moves from rank \rankDinohpFrozen{} to \rankDinohpProbed{} in the opposite direction, and the
two $256\times128$ C-RADIOv4 configurations take the top two places probed, from ranks 6 and 3
frozen. The two orderings share no top-three membership.

This is not a story about small differences being reordered by noise. Frozen, SigLIP2-g leads
the best C-RADIOv4 configuration \mapSiglipgFrozen{} to \mapRadiohFrozen{}; probed, the same
comparison runs \mapRadiohProbed{} to \mapSiglipgProbed{} the other way. Both gaps are two
orders of magnitude above the seed noise of \cref{sec:setup}. The rankings genuinely disagree,
and each is a correct answer to a different question.

\subsection{Why contrastive encoders lose their lead}
The \emph{gain} column is where the mechanism is visible. A fitted head multiplies frozen
Market-1501 mAP by \gainDinohp{}$\times$ for DINOv3-H+ and \gainRadioh{}$\times$ for
C-RADIOv4-H, but by only \gainSiglipg{}$\times$ for SigLIP2-g and \gainSiglipso{}$\times$ for
SigLIP2-SO400M. SigLIP2 is the outlier at both scales, by roughly a factor of two, and the two
scales agree with each other, which argues against an idiosyncrasy of one checkpoint.

The explanation is the pre-training objective. A sigmoid or softmax contrastive loss over
image--text pairs optimises cosine similarity in the embedding space directly; the frozen space
of such a model is therefore already close to the geometry that zero-shot cosine retrieval
reads, and an affine map has correspondingly little to recover. DINOv3's self-supervised
objective and C-RADIOv4's distillation objective optimise no such thing: their raw cosine
similarity is poor, but the identity information is present and a linear readout finds it. The
practical consequence generalises past the encoders in this table:

\begin{quote}
\textbf{Frozen-cosine benchmarking systematically flatters contrastively pre-trained encoders.}
A backbone comparison conducted that way is measuring readout alignment as much as
representation quality, and will name a different winner than the same comparison conducted
with a fitted head.
\end{quote}

Had this study reported only the frozen column --- as the zero-shot convention would --- it
would have concluded that a language-aligned encoder is the best available frozen ReID
representation. Had it reported only the probed column, it would have concluded the opposite.
Both conclusions are supported by the data; neither is supported without naming the readout.

\subsection{What this does not claim}
It does not claim that a fitted head is the correct readout and a frozen one is the incorrect
one. Which readout is right is a question about the deployment: if no labelled data exists in
the target domain, the frozen column is the one that predicts what will happen, and SigLIP2-g's
lead there is real. The claim is narrower and, we think, harder to argue with --- that the
readout is an experimental axis with an effect size larger than the backbone differences the
experiment is usually run to detect, and that a comparison which fixes it silently is
under-specified.

\section{What a Head Buys, and Where the Gain Comes From}\label{sec:heads}
\Cref{sec:readout} established that the readout reorders the encoders. This section asks what
the readout is actually supplying, using the PCA control to separate two hypotheses that the
supervised rows cannot distinguish on their own.

\begin{table*}[t]
\centering
\caption{The lead encoder (C-RADIOv4-H at $224\times224$) under all four \texttt{market1501/train}-fitted readouts, mAP. Only
the Market-1501 column is in-domain; every other column applies the same Market-fitted head to
a domain whose labels it has never seen. The PCA row is the matched control --- same images,
same output dimension, no labels.}
\label{tab:heads}
\scriptsize
\setlength{\tabcolsep}{3pt}
% generated by tools/gen_tables.py from results/runs — edit the runs, not this
\begin{tabular}{lrrrrrrr}
\toprule
head & Market & MSMT17 & CUHK03 & Occ-REID & CCVID & VRIC & VRAI \\
\midrule
frozen & 0.061 & 0.023 & 0.016 & 0.456 & 0.523 & 0.119 & 0.174 \\
PCA & 0.069 & 0.027 & 0.021 & 0.456 & 0.573 & 0.106 & 0.163 \\
linear & 0.400 & 0.093 & 0.130 & 0.624 & \textbf{0.624} & 0.130 & 0.171 \\
ArcFace & \textbf{0.709} & \textbf{0.165} & \textbf{0.153} & \textbf{0.646} & 0.609 & \textbf{0.191} & \textbf{0.228} \\
\bottomrule
\end{tabular}

\end{table*}

\subsection{The gain is the labels, not the bottleneck}
Every fitted head in this study produces a 512-dimensional embedding from a
higher-dimensional frozen one, so any improvement it shows is confounded between two causes:
the supervision it received, and the dimensionality reduction itself. Without a control, ``the
ArcFace head beats the frozen encoder'' is indistinguishable from ``512 dimensions are enough''.

The PCA row is that control. It is fitted on the same \nHeadImages{} images, produces the same
512-dimensional embedding, and uses none of the \nHeadIdentities{} identity labels. On
the lead encoder it moves Market-1501 from \leadMarketFrozen{} to \leadMarketPCA{}, moves
Occluded-REID from \leadOccludedFrozen{} to \leadOccludedPCA{} --- a change well inside that
dataset's seed noise --- and \emph{loses} on VRAI, \leadVraiFrozen{} to \leadVraiPCA{}. The
bottleneck is worth approximately nothing. The supervised heads on the same features reach
\leadMarketLinear{} and \leadMarketArc{}, so the entire effect reported in \cref{sec:readout}
is attributable to the identity supervision.

\begin{table}[t]
\centering
\caption{The same recipe fitted on each of two training splits, read on both of them, on the
lead encoder. The diagonal is in-domain. Every off-diagonal cell is the same recipe, the same
encoder, the same fitting budget and the same test set, and differs only in where the head was
fitted --- which is what makes this a control rather than a comparison.}
\label{tab:fitsplit}
\small
% generated by tools/gen_tables.py from results/runs — edit the runs, not this
\begin{tabular}{llrr}
\toprule
readout & fitted on & Market & MSMT17 \\
\midrule
frozen & --- & 0.061 & 0.023 \\
\midrule
PCA & Market-1501 & \textbf{0.069} & 0.027 \\
PCA & MSMT17 & 0.063 & \textbf{0.026} \\
\midrule
linear & Market-1501 & \textbf{0.400} & 0.093 \\
linear & MSMT17 & 0.248 & \textbf{0.260} \\
\midrule
ArcFace & Market-1501 & \textbf{0.709} & 0.165 \\
ArcFace & MSMT17 & 0.378 & \textbf{0.488} \\
\bottomrule
\end{tabular}

\end{table}

\subsection{Half of what a head buys is where it was fitted}
\Cref{sec:setup} fixes every head in the preceding tables on \texttt{market1501/train}, which
makes one column of each of them in-domain and every other column a transfer number. That
convention is the usual one, and it is also a confound: it cannot separate \emph{a head helps}
from \emph{a head fitted here helps}, because it never varies the second clause.

\Cref{tab:fitsplit} varies it. The same three recipes, the same encoders, the same budget and
the same seed are refitted on \texttt{msmt17/train} --- \nHeadImagesMsmt{} images of
\nHeadIdentitiesMsmt{} identities --- and each is read on both training domains. Nothing else
changes, and no backbone is re-run: the frozen features of both training splits were already
cached, so the control costs a head fit and not a single forward pass.

The result is a mirror. Averaged over every encoder, ArcFace reads \fitArcfaceMarketHere{} mAP
on Market-1501 when it was fitted there and \fitArcfaceMarketElse{} when the identical recipe
was fitted on MSMT17; on MSMT17 the same two heads read \fitArcfaceMsmtHere{} and
\fitArcfaceMsmtElse{}. The swing is \fitSwingArcfaceMarket{} in one direction and
\fitSwingArcfaceMsmt{} in the other. A gap visible in one direction only could be a property of
MSMT17; one that appears in both directions at the same size is a property of fitting.

It is monotone in supervision, and the PCA control is again what makes that legible. The same
swap moves PCA by \fitSwingPcaMarket{}, the linear head by \fitSwingLinearMarket{} and ArcFace
by \fitSwingArcfaceMarket{}. An unsupervised rotation of the same features barely encodes which
domain it was fitted on; an angular-margin classifier over identity labels encodes it most.

This sharpens \cref{sec:readout}'s claim rather than weakening it. A head that has never seen a
Market-1501 label still lifts Market-1501 mAP from \fitFrozenMarket{} to
\fitArcfaceMarketElse{}, so most of what the readout supplies is a transferable metric geometry
and not memorised identities; fitting it in-domain then roughly doubles that again. Of the
total gain over the frozen encoder, \fitShareArcfaceMarket{} survives being fitted elsewhere on
Market-1501 and \fitShareArcfaceMsmt{} does on MSMT17. The readout still determines the
ranking, as \cref{sec:readout} reports. What the fit domain determines is how much of the
readout a deployment somewhere else gets to keep.

One asymmetry in \cref{tab:fitsplit} carries less than it appears to. \texttt{msmt17/train} is
\nHeadImagesMsmt{} images against Market-1501's \nHeadImages{}, so which of the two splits
yields the better \emph{transfer} head is confounded with how much data each head saw. The
mirror above is read within a column and does not rest on that comparison; a matched-budget
version of it is left to future work.

\subsection{A frozen backbone plus one affine layer is a real system}
The best in-domain configuration in this study reaches \bestMarketMAP{} mAP and
\bestMarketRone{} Rank-1 on \texttt{market1501/official@1}. This is the first number in the
study that may be placed beside a published one, because it is produced the way published
numbers are: fitted on the official training split, evaluated on the official protocol, with
the test identities disjoint from the training identities. It does not beat a tuned
specialist. It costs one forward pass over \nHeadImages{} cached crops plus well under a
minute of head fitting on a single consumer GPU, with no backbone gradient at any point.

\subsection{Angular margin, and where it fails}
ArcFace beats a plain linear head almost everywhere in this study, and by the largest margin
where the head was fitted --- \leadMarketLinear{} to \leadMarketArc{} on Market-1501. This is
expected: an angular margin optimises the cosine geometry the retrieval metric reads, and a
softmax cross-entropy head does not. The exceptions are informative and are concentrated in one
encoder family; \cref{tab:margin} shows them, and \cref{sec:ablation} interprets them.

\begin{table*}[t]
\centering
\caption{Relative change in mAP from replacing the linear head with the ArcFace head, same
features and same budget. Positive is the expected direction. The losses cluster in DINOv3's
feature space on transfer sets, and both DINOv3 probes reach a perfect training accuracy, so
they are not fit failures.}
\label{tab:margin}
\scriptsize
\setlength{\tabcolsep}{3pt}
% generated by tools/gen_tables.py from results/runs — edit the runs, not this
\begin{tabular}{lrrrrrrr}
\toprule
encoder & Market & MSMT17 & CUHK03 & Occ-REID & CCVID & VRIC & VRAI \\
\midrule
C-RADIOv4-H & +77\% & +76\% & +17\% & +3\% & \loss{-2\%} & +47\% & +33\% \\
C-RADIOv4-SO400M & +89\% & +77\% & +36\% & +7\% & +3\% & +41\% & +41\% \\
SigLIP2-g & +51\% & +28\% & +5\% & +3\% & +1\% & +39\% & +38\% \\
SigLIP2-SO400M & +78\% & +32\% & +16\% & +5\% & +2\% & +56\% & +36\% \\
DINOv3-H+ & +105\% & +70\% & \loss{-20\%} & \loss{-3\%} & \loss{-4\%} & +10\% & +32\% \\
DINOv3-L & +101\% & +53\% & +1\% & \loss{-9\%} & +4\% & +2\% & +17\% \\
CLIP-B/16 & +103\% & +4\% & \loss{-19\%} & +0\% & +14\% & +25\% & +11\% \\
\bottomrule
\end{tabular}

\end{table*}

\subsection{The head transfers out of its object class}
The same head fitted on \nHeadIdentities{} pedestrian identities improves aerial \emph{vehicle}
retrieval, from \leadVraiFrozen{} frozen to \leadVraiArc{} with ArcFace, having never seen a
vehicle label. Since the matched PCA control \emph{falls} on that dataset, this is not a
dimensionality effect either. What the head learned is not ``what a person looks like'' but a
metric geometry that instance discrimination reuses across object classes --- which is the
most useful practical finding in this section and the one that most needs a larger transfer
grid before it is leaned on.

\section{Does Agglomeration Cost the Instance Margin?}\label{sec:ablation}
Agglomerative backbones distil several heterogeneous teachers into one student
\cite{ranzinger2024amradio,heinrich2025radiov25}. The concern this raises for
re-identification is specific: classification asks whether a \emph{class} is separable, ReID
asks whether one \emph{individual} is separable from a visually near-identical other, and
averaging teacher objectives is a plausible way to lose exactly that margin while every
classification benchmark continues to look healthy.

The concern is testable here rather than by analogy, because C-RADIOv4's teachers are in the
encoder set. \Cref{tab:ablation} is the ablation: the student at two capacities, its literal
teacher SigLIP2-g at 1163M, a capacity-matched SigLIP2 at 428M, and the DINOv3
family bracketing the student from above and below. All rows use the ArcFace head, under which
every probe in the study converged, and all teacher rows use the \texttt{squash} geometry so
that \cref{sec:geometry}'s confound is not silently inside this table.

\begin{table*}[t]
\centering
\caption{The teacher ablation, ArcFace mAP. C-RADIOv4 is a student of the SigLIP2 and DINOv3
families; SigLIP2-g is the literal teacher. \emph{role} names that relation. Market-1501 is the
only column whose identities any head has seen. DINOv3 stands in for a teacher family, not for
the actual DINOv3-7B teacher, which does not fit on this study's GPU.}
\label{tab:ablation}
\small
% generated by tools/gen_tables.py from results/runs — edit the runs, not this
\begin{tabular}{lllrrrrrrr}
\toprule
encoder & M & role & Market & MSMT17 & CUHK03 & Occ-REID & CCVID & VRIC & VRAI \\
\midrule
C-RADIOv4-H & 653 & student & \textbf{0.709} & 0.165 & 0.153 & 0.646 & 0.609 & 0.191 & 0.228 \\
C-RADIOv4-SO400M & 431 & student & 0.696 & 0.162 & 0.140 & \textbf{0.651} & 0.625 & 0.197 & 0.207 \\
\midrule
SigLIP2-g & 1163 & teacher & 0.608 & \textbf{0.231} & \textbf{0.283} & 0.649 & 0.626 & 0.115 & 0.229 \\
SigLIP2-SO400M & 428 & teacher & 0.514 & 0.200 & 0.174 & 0.558 & 0.589 & 0.075 & 0.177 \\
\midrule
DINOv3-H+ & 840 & teacher family & 0.655 & 0.067 & 0.097 & 0.356 & 0.461 & \textbf{0.359} & \textbf{0.247} \\
DINOv3-L & 303 & teacher family & 0.607 & 0.069 & 0.075 & 0.332 & 0.434 & 0.282 & 0.160 \\
\midrule
CLIP-B/16 & 87 & control & 0.289 & 0.027 & 0.017 & 0.414 & \textbf{0.795} & 0.035 & 0.039 \\
\bottomrule
\end{tabular}

\end{table*}

\subsection{In-domain, the margin improves rather than erodes}
On Market-1501 the student beats its literal teacher, \studentMarket{} against
\teacherMarket{} mAP, despite the teacher having 1.8$\times$ its parameters and running
at a higher input resolution that \cref{sec:geometry} shows is an advantage rather than a
neutral difference.

More to the point, the advantage is \emph{larger} on the metric the concern is about. mINP
\cite{ye2022deep} measures the rank of the hardest true match, which is the instance margin
that distillation was predicted to destroy. In-domain the student's mINP is
\leadMarketMINP{} against the teacher's \teacherMarketMINP{}, a factor of \minpGap{}$\times$
where the mAP factor is smaller. Given labels in the target domain, agglomerative distillation
does not merely preserve the fine-grained margin --- it improves it.

\subsection{Off-domain, the advantage does not survive}
It is a narrower claim than it first appears, and the rest of \cref{tab:ablation} says why.
Market-1501 is the only dataset whose identities any head has seen. On the five transfer sets
the student never exceeds the better teacher family by more than the seed noise of
\cref{sec:setup}: it ties on Occluded-REID (\familyStudentOccluded{} against
\familySiglipOccluded{}) and CCVID (\familyStudentCcvid{} against \familySiglipCcvid{}), and
loses on CUHK03-NP, VRIC and VRAI. On CUHK03-NP it loses to a teacher 1.5$\times$
\emph{smaller} than itself.

So the property is not ``agglomeration wins on people'' --- CUHK03-NP is a person dataset and
falsifies that --- and it is not ``agglomeration wins'' either. It is that \textbf{the
student's gain is concentrated exactly where the probe had labels}, which is also why the
retention view of the same numbers (\cref{tab:retention}) is not independent evidence: high
in-domain performance with mediocre transfer \emph{is} low retention, and reporting both as
two findings would be counting one twice.

\begin{table*}[t]
\centering
\caption{Retention: target mAP divided by source mAP, for the one head fitted on
\texttt{market1501/train}. The absolute source column is printed beside the ratios on purpose
--- a retention ratio flatters a weak source model, and CLIP's Occluded-REID and CCVID entries
are exactly that trap.}
\label{tab:retention}
\scriptsize
\setlength{\tabcolsep}{3pt}
% generated by tools/gen_tables.py from results/runs — edit the runs, not this
\begin{tabular}{lrrrrrrr}
\toprule
encoder & Market (source) & MSMT17 & CUHK03 & Occ-REID & CCVID & VRIC & VRAI \\
\midrule
C-RADIOv4-H & 0.709 & 0.23 & 0.22 & 0.91 & 0.86 & 0.27 & 0.32 \\
C-RADIOv4-SO400M & 0.696 & 0.23 & 0.20 & 0.94 & 0.90 & 0.28 & 0.30 \\
SigLIP2-g & 0.608 & 0.38 & 0.46 & 1.07 & 1.03 & 0.19 & 0.38 \\
SigLIP2-SO400M & 0.514 & 0.39 & 0.34 & 1.09 & 1.15 & 0.15 & 0.35 \\
DINOv3-H+ & 0.655 & 0.10 & 0.15 & 0.54 & 0.70 & 0.55 & 0.38 \\
DINOv3-L & 0.607 & 0.11 & 0.12 & 0.55 & 0.71 & 0.46 & 0.26 \\
CLIP-B/16 & 0.289 & 0.09 & 0.06 & 1.43 & 2.75 & 0.12 & 0.14 \\
\bottomrule
\end{tabular}

\end{table*}

\subsection{The transfer loss has a shape, and it is predictable}
The interesting result is not that the student loses off-domain but \emph{where}. The two
teacher families disagree sharply and by domain: language-aligned features carry degraded
person crops, and self-supervised dense features carry vehicles. SigLIP2 leads DINOv3 by
\spreadCuhk{}$\times$ on CUHK03-NP and by \spreadMsmt{}$\times$ on MSMT17; DINOv3 leads SigLIP2
by \spreadVric{}$\times$ on VRIC. \Cref{tab:bracket} places the student inside that
disagreement.

\begin{table*}[t]
\centering
\caption{Best-of-family ArcFace mAP, with the teachers' spread and the student's position
inside it. \emph{spread} is the ratio of the better teacher family to the worse; \emph{vs.\
best} is the student's relative distance from the better one. The two datasets where the
teachers disagree by roughly 3$\times$ both cost the student roughly
45\%.}
\label{tab:bracket}
\scriptsize
\setlength{\tabcolsep}{3pt}
% generated by tools/gen_tables.py from results/runs — edit the runs, not this
\begin{tabular}{lrrrrcr}
\toprule
dataset & SigLIP2 & DINOv3 & C-RADIOv4 & spread & bracketed & vs.\ best \\
\midrule
Market & 0.608 & 0.655 & \textbf{0.709} & 1.08$\times$ & \loss{above} & +8.2\% \\
MSMT17 & 0.231 & 0.069 & \textbf{0.165} & 3.35$\times$ & yes & \loss{-28.7\%} \\
CUHK03 & 0.283 & 0.097 & \textbf{0.153} & 2.91$\times$ & yes & \loss{-45.9\%} \\
Occ-REID & 0.649 & 0.356 & \textbf{0.651} & 1.82$\times$ & \loss{above} & +0.2\% \\
CCVID & 0.626 & 0.461 & \textbf{0.625} & 1.36$\times$ & yes & \loss{-0.0\%} \\
VRIC & 0.115 & 0.359 & \textbf{0.197} & 3.12$\times$ & yes & \loss{-45.1\%} \\
VRAI & 0.229 & 0.247 & \textbf{0.228} & 1.08$\times$ & \loss{below} & \loss{-7.7\%} \\
\bottomrule
\end{tabular}

\end{table*}

The pattern is regular enough to state as a rule of thumb. Where the teacher families agree
--- Occluded-REID at \spreadOccluded{}$\times$, CCVID at \spreadCcvid{}$\times$, VRAI at
\spreadVrai{}$\times$ --- the student matches the better of them to within
8\% and usually to within noise. Where they disagree by roughly a factor of three
--- CUHK03-NP at \spreadCuhk{}$\times$ and VRIC at \spreadVric{}$\times$ --- the student falls
\deficitCuhk{}\% and \deficitVric{}\% short of the better teacher respectively. Two independent
datasets, one person and one vehicle, at nearly the same teacher spread, produce nearly the
same student deficit.

\textbf{This is a different failure than the one the literature warns about, and a more useful
one.} The cost of agglomeration here is not a destroyed instance margin --- \cref{sec:ablation}
opens by showing the margin improving where labels exist. It is \emph{regression toward the
teacher mean on any domain where one teacher is much better than the other}. That cost is
predictable in advance from a question a practitioner can answer without running the
experiment: do my images look more like what the language-aligned teacher was good at, or more
like what the self-supervised one was good at? Where the honest answer is ``strongly one of
them'', the blend is the wrong purchase and the better single teacher is available.

We report one prospective check on this. Both C-RADIOv4 rows for CUHK03-NP were forecast to
fall inside the teacher bracket before they were run, on the strength of the pattern in the
datasets already measured, and both did.

\subsection{A family signature: the angular margin does not transfer for DINOv3}
\Cref{tab:margin} showed that replacing the linear head with an ArcFace head helps almost
everywhere. The exceptions are concentrated in one family: in DINOv3's feature space, on
transfer sets, the angular margin loses. Both DINOv3 probes reach a perfect training accuracy
over the source identities, so this is not an optimisation failure --- the head fits, and what
it fits transfers worse than the plainer head does. The practical reading is uncomfortable and
worth stating plainly: \emph{fit an angular-margin head on your labelled domain} is good advice
for an agglomerative or language-aligned backbone and bad advice for a DINOv3 backbone deployed
cross-domain, which is the opposite of what a Market-1501-only study would conclude.

\subsection{One anomaly we cannot yet explain}
CCVID is the one column in \cref{tab:ablation} where the ordering does not resemble any other.
The 87M CLIP control, run under the \texttt{squash} geometry, scores \clipCcvid{} mAP
there, ahead of every larger encoder in the study including its own 1163M
language-aligned relative. Its retention on that dataset is correspondingly implausible
(\cref{tab:retention}). The result is reproducible from the cached features and the protocol
digest is shared with every other row in the column, so it is not an adapter fault of the kind
\cref{sec:substrate}'s guards are designed to catch; but we do not have a mechanism for it, the
tracklet pooling step is unique to this dataset, and we flag it as unresolved rather than
interpret it. It is listed among the open items in \cref{sec:limitations}.

\section{Two Nuisance Variables Larger Than the Effect}\label{sec:geometry}
Neither of the variables in this section is interesting in itself. Both are reported because
each is larger than several of the backbone differences that frozen-encoder comparisons are run
to establish, and neither is routinely reported at all.

\subsection{How the crop is fitted to the input}
A pedestrian crop is roughly $64\times128$; a ViT input is square. Something has to reconcile
them, and the two conventions in common use disagree about what. \texttt{timm}'s default
evaluation transform resizes the short side and centre-crops, which on such an image retains
roughly the middle 45\% of the body and discards the head and the feet. The
\texttt{torchhub} path resizes to the target directly, distorting the aspect ratio and keeping
everything.

Measured on identical CLIP ViT-B/16 weights at identical input size, with geometry as the only
variable, the choice moves probed Market-1501 mAP from \clipcropMarket{} to \clipMarket{}. It
also moves the head's training accuracy over the source identities from $0.7073$ to $0.9319$,
which is the more diagnostic of the two: under the centre-crop convention an angular margin
cannot separate the training identities in CLIP's frozen space, and under the other convention
it nearly can. A conclusion of the form ``an angular margin fails in CLIP's representation''
--- which an earlier version of this study drew --- was substantially a statement about a
resize transform.

The consequence for the ablation of \cref{sec:ablation} is direct. The teachers reach the bench
through \texttt{timm} and the student through \texttt{torchhub}. Run on each backend's default,
the teachers would have carried a handicap the student never paid, worth up to a large fraction
of their probed score, and \cref{tab:ablation} would have reported the opposite conclusion with
the same confidence. This is why the resize geometry is a hashed field of the encoder
specification in our substrate rather than a constant in a runner script: the two geometries
address different cache entries and appear as different rows.

\subsection{Input resolution}
\Cref{tab:geometry} varies input size while holding the encoder, the head and the protocol
fixed. Three regularities hold across both C-RADIOv4 capacities.

\begin{table*}[t]
\centering
\caption{Input geometry, ArcFace mAP. Rows 1--3 and 4--6 vary resolution within one
checkpoint; rows 7--8 vary only the resize convention. \emph{native} feeds each image at its
own size, snapped to the patch grid and capped at 512 pixels.}
\label{tab:geometry}
\scriptsize
\setlength{\tabcolsep}{3pt}
% generated by tools/gen_tables.py from results/runs — edit the runs, not this
\begin{tabular}{lllrrrrrrr}
\toprule
encoder & input & fit & Market & MSMT17 & CUHK03 & Occ-REID & CCVID & VRIC & VRAI \\
\midrule
C-RADIOv4-H & 224$\times$224 & squash & 0.709 & 0.165 & 0.153 & 0.646 & 0.609 & 0.191 & 0.228 \\
C-RADIOv4-H & 256$\times$128 & squash & \textbf{0.718} & \textbf{0.180} & \textbf{0.163} & 0.622 & 0.612 & 0.185 & 0.187 \\
C-RADIOv4-H & \textit{native} & squash & 0.559 & 0.127 & 0.119 & 0.460 & 0.589 & 0.106 & \textbf{0.244} \\
\midrule
C-RADIOv4-SO400M & 224$\times$224 & squash & 0.696 & 0.162 & 0.140 & \textbf{0.651} & 0.625 & \textbf{0.197} & 0.207 \\
C-RADIOv4-SO400M & 256$\times$128 & squash & 0.718 & 0.175 & 0.149 & 0.647 & 0.581 & 0.177 & 0.179 \\
C-RADIOv4-SO400M & \textit{native} & squash & 0.573 & 0.136 & 0.119 & 0.481 & 0.612 & 0.097 & 0.225 \\
\midrule
CLIP-B/16 (squash) & 224$\times$224 & squash & 0.289 & 0.027 & 0.017 & 0.414 & \textbf{0.795} & 0.035 & 0.039 \\
CLIP-B/16 (crop) & 224$\times$224 & crop & 0.174 & 0.016 & 0.018 & 0.437 & 0.630 & 0.049 & 0.035 \\
\bottomrule
\end{tabular}

\end{table*}

First, \textbf{$256\times128$ wins on Market-1501 and CUHK03-NP and loses on both vehicle
sets}, at both C-RADIOv4 capacities, while also running about 1.6$\times$ faster; where it
wins it is therefore the better purchase on both axes at once.

Second, the tempting explanation for that --- that a $2{:}1$ input matches a person crop's
aspect ratio --- is \textbf{false}, and this study contains its own counterexample.
Market-1501 and Occluded-REID ship images that are every one of them exactly $64\times128$:
identical median size, identical $2.00$ aspect, coefficient of variation $0.00$ in image area.
With the same head they prefer \emph{opposite} input sizes --- $256\times128$ on Market
($0.718$ against $0.709$) and $224\times224$ on Occluded-REID ($0.646$ against $0.622$, about
five times that set's seed noise). CUHK03-NP, at a much taller $2.97$ aspect, prefers $2{:}1$
by \emph{more} than Market does rather than less. A second candidate explanation, token count,
fails on the same rows. We report the regularity and explicitly do not claim a mechanism for
it; ``$2{:}1$ usually wins on person crops'' is a useful default with two measured exceptions
inside this table, not a law.

Third, \textbf{\emph{native} resolution wins on exactly one dataset, VRAI, and loses
everywhere else --- including on the other vehicle set}. That last detail is the one that
identifies the variable. If the split were person-versus-vehicle, native should also win on
VRIC; instead it collapses there, from $0.191$ to $0.106$. What separates the two is crop
size. VRAI's median crop is $295\times202$, inside the range C-RADIOv4 was trained across;
VRIC's crops, and every Market-1501 and Occluded-REID crop, are far below it. Indeed on the two
person sets where every file is exactly $64\times128$, ``native'' is not native at all: it is a
fixed 32-token resolution wearing a misleading label, well under the roughly 128-pixel floor
the encoder saw in training.

The variable that moved is therefore resolution \emph{relative to the encoder's training
distribution}, not resampling and not object class. Upsampling a small crop to $224$ is not a
distortion to be avoided; it is how the crop reaches a size the encoder can read. Native is
also far cheaper on small crops --- roughly 4$\times$ the throughput --- so it remains a real
option wherever a quarter of the mAP is an acceptable price.

Two further cells resist the story and we report them rather than smooth them: CLIP's
Occluded-REID and VRIC rows prefer the centre-crop geometry, which every other column in the
table penalises. Both are among the weakest scores in the study, and we have not established
whether that is the reason.

\subsection{Why this belongs in the paper rather than in an appendix}
Take the two variables together. Resize geometry is worth up to a large fraction of a probed
score; resolution is worth on the order of a fifth to a quarter of one. Several of the backbone
comparisons in \cref{tab:ablation} are decided by less than either. A frozen-encoder comparison
that does not pin both is therefore not merely imprecise --- it is capable of reporting the
reverse of the truth, and it will do so with a table that looks exactly like a correct one.
The remedy is cheap and we recommend it as a reporting standard: state the resize convention
and the input size beside every frozen-encoder number, and treat any comparison that omits
them as under-specified.

\begin{table}[t]
\centering
\caption{Extraction throughput on this study's single RTX 2070 Max-Q, Market-1501, frozen
features. A throughput figure measures a machine and a moment, not a model --- one cell here
first recorded 2.2 img/s against 12.8 for the same encoder elsewhere because
another process held part of the card; re-extracting on a quiet card gave 12.3, and the
metrics never moved.}
\label{tab:cost}
\scriptsize
% generated by tools/gen_tables.py from results/runs — edit the runs, not this
\begin{tabular}{lllr}
\toprule
encoder & input & params (M) & img/s \\
\midrule
C-RADIOv4-H & 224$\times$224 & 653 & 12.3 \\
C-RADIOv4-H & 256$\times$128 & 653 & 20.8 \\
C-RADIOv4-H & \textit{native} & 653 & 53.9 \\
C-RADIOv4-SO400M & 224$\times$224 & 431 & 21.3 \\
C-RADIOv4-SO400M & 256$\times$128 & 431 & 31.4 \\
C-RADIOv4-SO400M & \textit{native} & 431 & 85.9 \\
SigLIP2-g & 256$\times$256 & 1163 & 5.8 \\
SigLIP2-SO400M & 224$\times$224 & 428 & 15.6 \\
DINOv3-H+ & 224$\times$224 & 840 & 9.5 \\
DINOv3-L & 224$\times$224 & 303 & 25.5 \\
CLIP-B/16 (squash) & 224$\times$224 & 87 & 81.8 \\
CLIP-B/16 (crop) & 224$\times$224 & 87 & 33.9 \\
\bottomrule
\end{tabular}

\end{table}

\section{From Rankings to Decisions}\label{sec:decisions}
% ---------------------------------------------------------------------------------------
% SCAFFOLD. This section is planned, not measured. The substrate already implements the
% metrics named below (reidbench.measure.openset, .calibration, .selective, .cascade) and
% the open-set protocol; what is missing is the sweep and the rejection head.
%
% Everything above this line is measured and citable today. Everything below is an outline
% with the claims it intends to make, so that the experiments can be dropped in without
% restructuring the paper. Delete this comment and the \scaffold markers when the runs land.
% ---------------------------------------------------------------------------------------
Every number reported so far answers a ranking question: given that the right answer is
somewhere in the gallery, how highly is it ranked? Deployments do not ask that question. They
ask whether a candidate match should be accepted at all, at a threshold fixed before the query
arrives, against a gallery that frequently does not contain the answer. A confident wrong match
in that setting is not a lower mAP --- it is the wrong person.

This section reports what the encoders and heads of \cref{sec:setup} do when they are scored as
\emph{decisions} rather than as rankings. The substrate of \cref{sec:substrate} already
expresses these measures --- the same $(\mathrm{rel}, \mathrm{valid})$ pair supports them
without a new abstraction --- so the additional cost is the sweep, not new machinery.

\subsection{Open-set protocol}
\scaffold{To be measured. A protocol in which a declared fraction of queries have no true match
in the gallery, with the non-mate set drawn under a fixed, published rule and stratified into
near and far difficulty following OpenOOD's discipline \cite{zhang2023openood}. Metrics:
DIR@FAR and FNIR@FPIR after the biometric identification convention
\cite{liao2014operid}, AUROC for the mated-versus-non-mated decision, and equal error rate.
The claim under test: that the ranking of \cref{tab:readout} and the ranking by decision
quality disagree, and specifically that the readout which wins on mAP is not the readout which
wins on FNIR at a usable FPIR.}

\subsection{Calibration}
\scaffold{To be measured. Expected calibration error and reliability curves for the similarity
score read as a confidence, per encoder and per head, in-domain and on each transfer set. The
claim under test: that cosine similarity from a margin-trained head is systematically
overconfident off its training domain, and that the overconfidence is worse for the encoder
that wins in-domain --- i.e.\ that \cref{sec:ablation}'s in-domain/transfer asymmetry has a
calibration signature as well as an accuracy one.}

\subsection{A calibrated rejection head}
\scaffold{To be built and measured. A distance-based readout with a parameter-free abstain
option, fitted on the same cached frozen features and costing nothing additional at inference
\cite{halo2026,dhamija2018agnostophobia}. Compared against the ArcFace head of
\cref{sec:setup} at matched embedding dimension and matched fitting budget, on both the ranking
metrics of \cref{sec:readout} and the decision metrics above. The claim under test: that the
abstain option can be added without cost to closed-set retrieval quality, which is what would
make it adoptable.}

\subsection{Threshold transfer}
\scaffold{To be measured. An acceptance threshold selected on the source domain, held fixed,
and applied to each transfer domain; reported as the realised operating point against the
intended one. This is the quantity a deployment actually needs and the one no benchmark in
\cref{sec:related} reports. The claim under test: that thresholds selected on Market-1501 do
not hold on any transfer domain, and that the size of the miss is predicted by the calibration
measurement above rather than by the mAP gap.}

\subsection{Why these belong with the preceding sections}
The connection is not thematic. \Cref{sec:readout} showed that the readout determines the
ranking; the readout also determines the \emph{geometry of the score}, and a threshold is a cut
through that geometry. An encoder chosen by frozen cosine similarity, an encoder chosen by
probed mAP, and an encoder chosen by decision quality at a fixed operating point need not be
the same encoder, and the first two are already known here to differ. Establishing whether the
third differs from both is the natural completion of this study, and the substrate that
produced \cref{tab:readout} is the one that will produce it.

\section{Limitations and Threats to Validity}\label{sec:limitations}
We state these at length because several of them bound the conclusions above rather than
merely qualifying them.

\paragraph{One head-training domain}
Every fitted head in this study is fitted on \texttt{market1501/train}. That is what makes the
in-domain versus transfer contrast clean, and it is also the study's largest structural limit:
``the student's gain is concentrated where the probe had labels'' is established with one
source domain. Whether the same asymmetry appears with a head fitted on MSMT17 or on a vehicle
set is not tested here, and it is the first experiment we would add.

\paragraph{MSMT17 is measured for the teachers only}
\Cref{tab:coverage} shows the gap. MSMT17 is the large-scale person transfer target and the
dataset best placed to settle the retention question of \cref{tab:retention} on a
gallery-comparable footing, and the C-RADIOv4 rows for it are not yet run. The retention ratios
we do report are not gallery-matched --- CUHK03-NP's gallery is roughly a third of
Market-1501's --- and should not be placed beside published retention figures without that
caveat.

\paragraph{DINOv3 is a stand-in}
C-RADIOv4's actual DINOv3 teacher is the 7B model, which does not fit on the GPU this
study ran on. The DINOv3 rows bracket the student's capacity from above and below and represent
the teacher \emph{family}; SigLIP2-g is the genuine article. Conclusions about the DINOv3 side
of the bracket are correspondingly weaker than those about the SigLIP2 side.

\paragraph{The CCVID anomaly}
CLIP ViT-B/16 at 87M scores \clipCcvid{} mAP on CCVID, ahead of every larger encoder
here. We have verified that it is not an adapter fault --- the protocol and manifest digests are
shared with every other row in that column --- and we have no mechanism for it. CCVID is the
one dataset in the study whose features pass through a tracklet pooling step, which is the
obvious place to look and which we have not yet ruled in or out. Until it is explained, we do
not draw any conclusion from the CCVID column, and readers should not either.

\paragraph{Two protocols are non-standard, for reasons outside our control}
Occluded-REID ships no standard split and labels no cameras, so its protocol uses the whole set
and cannot apply the same-camera exclusion that every Market-1501 number depends on. VRAI
withholds its test identities, so the only offline-scorable protocol runs over its training
split. That is sound for every row in this paper --- no head here has seen a VRAI image --- and
it would be unsound for any model fine-tuned on VRAI, which is a constraint on who may cite
these rows rather than on the rows themselves.

\paragraph{Absolute values are not comparable across datasets}
Gallery sizes here span roughly a thousand to over thirty thousand images. Every claim we make
is within a column, where all rows share a protocol digest and a manifest digest. The tables
print query and gallery sizes (\cref{tab:datasets}) so that any reader tempted to compare
across columns can see the reason not to.

\paragraph{Throughput measures a machine}
The figures in \cref{tab:cost} come from one thermally throttled consumer GPU and are useful
only as ratios within that table, as the caption's example shows.

\paragraph{Licence heterogeneity}
Of the checkpoints here, only the two C-RADIOv4 weights carry a verified licence permitting
commercial use; the rest are unverified or research-only, as are most of the datasets. Our
substrate reports this per run. It does not change any measurement and it changes the practical
recommendation considerably.

\paragraph{No backbone was trained}
This is a design choice, not an oversight, and it bounds the claims accordingly. Nothing here
speaks to what these encoders do when fine-tuned, and \cref{sec:readout}'s reordering result
should not be assumed to survive fine-tuning --- probe rankings are known not to predict
fine-tuned rankings in general, and testing that here would have cost more compute than the
entire study.

\section{Conclusion}
We evaluated \nEncoders{} frozen vision foundation backbones at \nEncoderConfigs{}
encoder-resolution configurations, under \nHeads{} readouts, across \nDatasets{} person and
vehicle re-identification datasets and \nProtocols{} exact protocols, for \nRuns{} run records
in total, without training a single backbone.

The result we would most like to see acted on is methodological. \textbf{The readout is an
experimental axis, and it is a larger one than the backbone differences that frozen-encoder
comparisons exist to detect.} On Market-1501, changing only how cached features are read moves
SigLIP2-g from rank \rankSiglipgFrozen{} to rank \rankSiglipgProbed{}, and DINOv3-H+ from rank
\rankDinohpFrozen{} to rank \rankDinohpProbed{}. Contrastively pre-trained encoders gain least from a fitted head because
their objective already optimised the geometry that zero-shot cosine retrieval reads, so
frozen-cosine benchmarking flatters them systematically. A matched label-free control shows the
gain is the supervision and not the 512-dimensional bottleneck. Any comparison of frozen
encoders that does not name its readout --- along with its resize convention and input size,
each of which we show to be worth more than several of the differences being reported --- is
under-specified, and can report the reverse of the truth with a table that looks correct.

On the specific question of agglomerative distillation, the answer is two-sided and neither
side is the one the concern anticipated. Where the probe has labels, distillation does not cost
the instance margin --- it improves it, by a wider factor on mINP (\minpGap{}$\times$) than on
mAP. Where the probe does not have labels, the student regresses toward the mean of teachers
that disagree sharply by domain, and the size of that regression tracks the size of the
disagreement: two independent datasets at a teacher spread near 3$\times$ both cost the
student close to 45\% relative. That makes the cost of agglomeration predictable
from a question a practitioner can answer without running the experiment. The practical
recommendation follows and needs its qualifier attached: an agglomerative backbone is the right
default \emph{if you can fit a head on labelled data from your own domain}; if you cannot, the
better single teacher for your domain is available and is better.

Finally, a frozen agglomerative backbone with one 512-dimensional affine head reaches
\bestMarketMAP{} mAP and \bestMarketRone{} Rank-1 on \texttt{market1501/official@1}, for one
forward pass over cached crops plus well under a minute of fitting on a consumer GPU. It does
not beat a tuned specialist. It is a reminder of how much of a modern ReID number is already
present in a general-purpose representation, waiting on the readout.

The natural completion of this work is \cref{sec:decisions}: to ask the same encoders and heads
not how well they rank, but how well they decide --- whether a match should be accepted at all,
at a threshold fixed in advance, against a gallery that may not contain the answer. The
substrate that produced every table here already expresses those measures.

\section*{Reproducibility}
The evaluation package, the protocol definitions with their digests, the encoder and head
specifications, and the complete run records --- including per-query results, the provenance of
every row, and the licence index --- are released under the MIT licence. The tables in this
paper are generated from those records by a script included with the source; no metric in this
document was transcribed by hand. We redistribute no dataset; obtaining each one, and complying
with its licence, remains the reader's responsibility.

\appendices
\section{Complete Result Matrices}\label{sec:appendix}
The tables in the body are slices chosen to make one comparison each. These three are the full
grid, so that any comparison the body does not make can still be checked. All rows within a
column share a protocol digest and a manifest digest; no row is comparable across columns
(\cref{sec:setup}). CCVID is shown under \texttt{ccvid/tracklet@1}; its cloth-changing
protocol is a separate question and is reported separately in the released records.

\begin{table*}[htbp]
\centering
\caption{Frozen encoders, ranked by cosine similarity over the raw pooled features, mAP. This
is the zero-shot convention, and it is the ordering \cref{sec:readout} shows is not stable.}
\label{tab:appendix-frozen}
% generated by tools/gen_tables.py from results/runs — edit the runs, not this
\begin{tabular}{llrrrrrrr}
\toprule
encoder & input & Market & MSMT17 & CUHK03 & Occ-REID & CCVID & VRIC & VRAI \\
\midrule
C-RADIOv4-H & 224$\times$224 & 0.061 & 0.023 & 0.016 & 0.456 & 0.523 & 0.119 & 0.174 \\
C-RADIOv4-H & 256$\times$128 & 0.059 & 0.024 & 0.019 & 0.447 & 0.523 & 0.107 & 0.138 \\
C-RADIOv4-H & \textit{native} & 0.044 & 0.019 & 0.017 & 0.338 & 0.509 & 0.074 & 0.194 \\
C-RADIOv4-SO400M & 224$\times$224 & 0.063 & 0.024 & 0.016 & 0.444 & 0.523 & 0.125 & 0.149 \\
C-RADIOv4-SO400M & 256$\times$128 & 0.063 & 0.024 & 0.016 & 0.440 & 0.536 & 0.111 & 0.127 \\
C-RADIOv4-SO400M & \textit{native} & 0.047 & 0.020 & 0.014 & 0.340 & 0.520 & 0.068 & 0.160 \\
SigLIP2-g & 256$\times$256 & \textbf{0.105} & \textbf{0.081} & \textbf{0.084} & \textbf{0.500} & 0.633 & 0.058 & 0.164 \\
SigLIP2-SO400M & 224$\times$224 & 0.077 & 0.076 & 0.044 & 0.403 & 0.574 & 0.036 & 0.127 \\
DINOv3-H+ & 224$\times$224 & 0.049 & 0.012 & 0.020 & 0.215 & 0.389 & \textbf{0.357} & \textbf{0.235} \\
DINOv3-L & 224$\times$224 & 0.051 & 0.016 & 0.015 & 0.282 & 0.351 & 0.299 & 0.149 \\
CLIP-B/16 (squash) & 224$\times$224 & 0.029 & 0.008 & 0.007 & 0.312 & \textbf{0.746} & 0.023 & 0.031 \\
CLIP-B/16 (crop) & 224$\times$224 & 0.023 & 0.006 & 0.007 & 0.280 & 0.561 & 0.027 & 0.025 \\
\bottomrule
\end{tabular}

\end{table*}

\begin{table*}[htbp]
\centering
\caption{The ArcFace head fitted on \texttt{market1501/train}, applied unchanged everywhere,
mAP. Market-1501 is the only in-domain column.}
\label{tab:appendix-map}
% generated by tools/gen_tables.py from results/runs — edit the runs, not this
\begin{tabular}{llrrrrrrr}
\toprule
encoder & input & Market & MSMT17 & CUHK03 & Occ-REID & CCVID & VRIC & VRAI \\
\midrule
C-RADIOv4-H & 224$\times$224 & 0.709 & 0.165 & 0.153 & 0.646 & 0.609 & 0.191 & 0.228 \\
C-RADIOv4-H & 256$\times$128 & \textbf{0.718} & 0.180 & 0.163 & 0.622 & 0.612 & 0.185 & 0.187 \\
C-RADIOv4-H & \textit{native} & 0.559 & 0.127 & 0.119 & 0.460 & 0.589 & 0.106 & 0.244 \\
C-RADIOv4-SO400M & 224$\times$224 & 0.696 & 0.162 & 0.140 & \textbf{0.651} & 0.625 & 0.197 & 0.207 \\
C-RADIOv4-SO400M & 256$\times$128 & 0.718 & 0.175 & 0.149 & 0.647 & 0.581 & 0.177 & 0.179 \\
C-RADIOv4-SO400M & \textit{native} & 0.573 & 0.136 & 0.119 & 0.481 & 0.612 & 0.097 & 0.225 \\
SigLIP2-g & 256$\times$256 & 0.608 & \textbf{0.231} & \textbf{0.283} & 0.649 & 0.626 & 0.115 & 0.229 \\
SigLIP2-SO400M & 224$\times$224 & 0.514 & 0.200 & 0.174 & 0.558 & 0.589 & 0.075 & 0.177 \\
DINOv3-H+ & 224$\times$224 & 0.655 & 0.067 & 0.097 & 0.356 & 0.461 & \textbf{0.359} & \textbf{0.247} \\
DINOv3-L & 224$\times$224 & 0.607 & 0.069 & 0.075 & 0.332 & 0.434 & 0.282 & 0.160 \\
CLIP-B/16 (squash) & 224$\times$224 & 0.289 & 0.027 & 0.017 & 0.414 & \textbf{0.795} & 0.035 & 0.039 \\
CLIP-B/16 (crop) & 224$\times$224 & 0.174 & 0.016 & 0.018 & 0.437 & 0.630 & 0.049 & 0.035 \\
\bottomrule
\end{tabular}

\end{table*}

\begin{table*}[htbp]
\centering
\caption{The same configurations as \cref{tab:appendix-map}, Rank-1. Reported because mAP and
Rank-1 do not always agree --- notably on CCVID, where several encoders separate on one and not
on the other.}
\label{tab:appendix-r1}
% generated by tools/gen_tables.py from results/runs — edit the runs, not this
\begin{tabular}{llrrrrrrr}
\toprule
encoder & input & Market & MSMT17 & CUHK03 & Occ-REID & CCVID & VRIC & VRAI \\
\midrule
C-RADIOv4-H & 224$\times$224 & 0.879 & 0.390 & 0.151 & 0.716 & 0.650 & 0.125 & 0.268 \\
C-RADIOv4-H & 256$\times$128 & \textbf{0.887} & 0.415 & 0.158 & 0.702 & 0.646 & 0.120 & 0.220 \\
C-RADIOv4-H & \textit{native} & 0.801 & 0.318 & 0.108 & 0.539 & 0.635 & 0.066 & 0.283 \\
C-RADIOv4-SO400M & 224$\times$224 & 0.872 & 0.383 & 0.135 & 0.706 & 0.658 & 0.129 & 0.244 \\
C-RADIOv4-SO400M & 256$\times$128 & 0.885 & 0.402 & 0.150 & 0.723 & 0.637 & 0.111 & 0.213 \\
C-RADIOv4-SO400M & \textit{native} & 0.807 & 0.339 & 0.116 & 0.567 & 0.682 & 0.054 & 0.266 \\
SigLIP2-g & 256$\times$256 & 0.826 & \textbf{0.529} & \textbf{0.313} & \textbf{0.749} & 0.643 & 0.073 & 0.271 \\
SigLIP2-SO400M & 224$\times$224 & 0.756 & 0.481 & 0.200 & 0.631 & 0.598 & 0.042 & 0.215 \\
DINOv3-H+ & 224$\times$224 & 0.858 & 0.200 & 0.085 & 0.401 & 0.650 & \textbf{0.263} & \textbf{0.307} \\
DINOv3-L & 224$\times$224 & 0.829 & 0.211 & 0.077 & 0.372 & 0.615 & 0.198 & 0.209 \\
CLIP-B/16 (squash) & 224$\times$224 & 0.539 & 0.105 & 0.016 & 0.520 & \textbf{0.838} & 0.017 & 0.052 \\
CLIP-B/16 (crop) & 224$\times$224 & 0.356 & 0.066 & 0.019 & 0.523 & 0.663 & 0.025 & 0.047 \\
\bottomrule
\end{tabular}

\end{table*}


\bibliographystyle{IEEEtran}
\bibliography{refs}

\end{document}
